\chapter{Method and Progress}\label{ch:method}

This chapter describes the implemented method and the progress completed by the
midterm stage. The project is not yet at the final result-analysis stage, but
the first 200-episode Vermont runs have been completed. The strongest outcomes
so far are implementation and design outcomes: the training pipeline is
operational, multiple baselines have been implemented, and the communication
variants have been integrated into a common MAPPO framework. The current
numerical results are treated as preliminary reference points.


\section{Problem Definition}

The CityLearn Annex 96 task is formulated as a cooperative partially observable
stochastic game. There are $N=25$ building agents. At time step $t$, building
$i$ observes a local vector $o_i^t$ and selects a continuous action $a_i^t$ for
its flexible devices. The joint action changes building states, battery states,
thermal conditions, and aggregate portfolio load. The portfolio should track a
daily reference load profile $r_t$ while avoiding excessive comfort violations
and negative secondary impacts.


The project objective is not only to maximize training reward. The practical
objective is to reduce testing-month tracking error, measured by NMBE and
CV-RMSE, while keeping comfort, ramping, peak demand, and fairness within
acceptable ranges. Cost and carbon emissions are part of the Annex 96 secondary
metric set, but the CE1 schemas used here do not provide pricing or
carbon-intensity time series, so those two metrics are reported as unavailable
in the current experiments. This objective motivates the use of a centralized
critic and structured actor communication.


\section{Environment and Baseline Pipeline}

The first completed component is the environment interface. CityLearn is wrapped
for several learning backends:


\begin{itemize}
  \item independent PPO and independent SAC baselines, where each building
  learns without communication;
  \item standard MAPPO, where all buildings use a shared policy and a
  centralized critic;
  \item grouped MAPPO, where clusters of buildings share group-specific actors;
  \item communication-aware grouped MAPPO, where a communication layer is
  inserted between actor encoders and action heads.
\end{itemize}

The MAPPO wrapper returns local observations, centralized observations,
continuous action spaces, rewards, done masks, and episode lengths in the tensor
format expected by the on-policy MAPPO implementation. Reusing the established
MAPPO optimizer reduces implementation risk and keeps the research focus on
CityLearn adaptation, grouping, and communication.


\section{Grouped Actor Architecture}

The grouped actor design is based on selective parameter sharing. Full sharing
is sample efficient but may be too restrictive for a heterogeneous portfolio.
Fully independent actors are expressive but require more data and weaken
coordination. The implemented grouped actor architecture balances these two
extremes.


Buildings are clustered using K-means over flexibility-related features. The
current implementation uses features such as battery energy storage capacity
and aggregate HVAC capacity. After clustering, buildings in group $G_k$ share
the same actor network. All groups share one centralized critic. The critic
observes the concatenated portfolio state, while each actor receives local
building observations plus any communication-enhanced hidden features.


This architecture has been implemented in the \texttt{mappo\_grouped} package.
The same grouping information is reused by communication-aware variants so that
communication effects can be compared against a consistent grouped baseline.


\section{MAPPO Learning Structure}

The learning structure follows centralized training and decentralized
execution. For each group, the actor maps observations to a continuous action
distribution. The centralized critic estimates the value of the portfolio state.
After each rollout episode, generalized advantage estimation is computed and
the PPO clipped objective updates the policies.


The grouped replay buffer preserves the agent dimension during PPO mini-batch
sampling. This detail is important: if group-agent samples are flattened too
early, communication layers cannot reconstruct the correct set of agents for
message passing. The implemented buffer therefore keeps tensors in a shape that
supports both MAPPO updates and communication over group members.


\section{Communication Modules}

The communication layer is inserted between the actor encoder and action head:
\[
o_i \rightarrow \mathrm{encoder}_k \rightarrow h_i
\rightarrow \mathrm{communication} \rightarrow \tilde{h}_i
\rightarrow \mathrm{action\ head}_k \rightarrow a_i .
\]
Because the communication module is registered as part of the actor, its
parameters are optimized by the same PPO actor loss.


\subsection{No Communication}

The identity module is used as the no-communication baseline:
\[
\tilde{h}_i = h_i .
\]
This isolates the effect of message passing from the effect of grouping.


\subsection{CommNet-Style Mean Communication}

For each group, the CommNet-style module computes the mean hidden feature of
the other agents:
\[
m_i = \frac{1}{|G_k|-1}\sum_{j \in G_k, j \ne i} h_j .
\]
The message is transformed by an MLP and added through a residual connection:
\[
\tilde{h}_i = h_i + f_\theta(m_i).
\]
This is the simplest differentiable communication baseline.


\subsection{PowerNet-Style Neighbor Communication}

The PowerNet-style module restricts communication to a graph. The implemented
topologies are ring, chain, and full graph. For adjacency matrix $A$, the
neighbor message is
\[
m_i = \frac{1}{\max(1,\sum_j A_{ij})}\sum_j A_{ij}g_\theta(h_j).
\]
This tests whether sparse local communication can preserve coordination while
reducing noisy all-to-all message passing.


\subsection{TarMAC-Style Attention}

The TarMAC-style module computes query, key, and value projections for all
agents. Attention weights are
\[
\alpha_{ij} =
\frac{\exp(q_i^\top k_j / \sqrt{d})}
{\sum_l \exp(q_i^\top k_l / \sqrt{d})},
\]
and the received context is
\[
c_i = \sum_j \alpha_{ij}v_j .
\]
The updated feature is computed from $h_i$ and $c_i$. This module tests whether
learned receiver selection is beneficial for CityLearn coordination.


\subsection{Weighted Same-Group and Other-Group Communication}

The project-specific communication module separates same-group and other-group
messages. For group $G_k$,
\[
m^{same}_k = \frac{1}{|G_k|}\sum_{j \in G_k} h_j,
\qquad
m^{other}_k = \frac{1}{N-|G_k|}\sum_{j \notin G_k} h_j .
\]
The combined message is
\[
m_k = \alpha m^{same}_k + \beta m^{other}_k,
\qquad \alpha > \beta \geq 0.
\]
The default implementation uses $\alpha=0.75$ and $\beta=0.25$. The purpose is
to encode a clear inductive bias: buildings with similar flexibility should
communicate more strongly, but portfolio-level tracking still requires some
global context.


\section{Completed Implementation}

At the midterm stage, the following implementation work has been completed:


\begin{itemize}
  \item independent PPO and SAC training scripts;
  \item standard MAPPO training and testing scripts;
  \item grouped MAPPO with K-means clustering;
  \item communication-aware grouped MAPPO with CommNet, PowerNet-style,
  TarMAC-style, DIAL-style, and weighted communication variants;
  \item shared centralized critic and shared value normalizer for grouped
  policies;
  \item grouped replay buffer preserving communication structure during PPO
  updates;
  \item checkpoint loading and saving for test-only evaluation;
  \item Annex 96 metric export, daily metric tables, and plotting utilities,
  with cost and carbon fields marked as unavailable when the CE1 data provides
  no pricing or carbon-intensity signals.
\end{itemize}

The most important conclusion from the completed work is that the project has a
common experimental interface for the main planned methods. The 200-episode
Vermont results also provide early evidence that communication changes the
tracking, comfort, and secondary-metric trade-offs. At this stage, global
CommNet is the most promising communication variant, but this is only a
preliminary observation.


\section{Current Limitations}

The current limitation is that the available numerical results are incomplete:
they cover Vermont only, use one seed, and are based on 200 training episodes.
Therefore, the report does not claim that one communication method is
definitively superior to another. The final comparison will require longer
training, Texas experiments, multiple seeds where possible, and a deeper
analysis of daily load trajectories. The CE1 datasets used in the current
pipeline also omit pricing and grid-carbon time series, so cost and carbon
secondary metrics cannot be interpreted without adding external assumptions.
Another current limitation is that K-means uses static flexibility features;
future work may replace or augment this with learned embeddings based on
building trajectories.

